\section{Field}
\label{sec:field}
A field assigns values to configurations of a system. The values considered here are matrix components whose magnitudes and relative directions determine their interactions. The physical question is how such interactions can produce the mass of a closed structure when the Yang--Mills Lagrangian contains no elementary gluon mass term.

\subsection{From field interaction to collective mass}
In the usual spacetime description, a gauge connection has matrix-valued components $\FieldYMPotential$, where $\FieldYMCoordinateIndex$ labels a spacetime direction. Using Hermitian matrices, equal to their conjugate transposes, its curvature, or field strength, is
\FieldYMStrengthDisplay
Here $\FieldYMCoupling$ is the coupling strength and
\FieldYMCommutatorDisplay
Directional derivatives measure changes of the connection, while the commutator records the order dependence of its internal transformations. This product term allows the field to interact with itself.

Squaring the field strength retains both propagation and interaction. With $c$ the speed of light and $\hbar$ the reduced Planck constant, choose units where $\FieldYMNaturalUnits$. The Yang--Mills Lagrangian density is
\FieldYMLagrangianDisplay
Repeated indices are contracted with the spacetime metric. Expanding the square gives quadratic, cubic and quartic terms in the connection, with no separate mass term proportional to $\FieldYMMassTerm$~\cite{yangMills1954}. For the gluon field the gauge group is $\ColorGroup$, and the linearized transverse waves are massless. Appendix~\ref{app:coupling} gives the weak-wave calculation.

Absence of an elementary mass term does not preclude collective invariant mass. Restoring $c$, a system with total energy $\FieldYMTotalEnergy$ and total momentum $\FieldYMTotalMomentum$ has invariant mass $\FieldYMCollectiveMass$ satisfying
\FieldYMInvariantMassDisplay
Even massless components can have a nonzero value of this invariant when their momenta cancel. That kinematic observation alone does not establish binding. A glueball is a gauge-invariant bound excitation of the interacting gluon field, with mass determined by its complete rest energy. The weak-wave comparison, a two-excitation example and a lattice result are given in Appendix~\ref{app:coupling}.

The construction follows field comparisons through coupled phase rotations to a closed geometry, then uses its stress and complete energy--momentum to obtain a mass expression. The mass statement retains the condition that the total four-momentum is defined, future-directed and timelike. A common positive threshold for all physical excitations is a distinct mass-gap question~\cite{jaffeWitten}. The first step here is to determine what information about the field components can generate the proposed geometry.

\subsection{Comparing components}
Magnitude alone does not specify an interaction between directions. Two components may have the same length and different alignments. An inner product measures length by self-pairing and alignment by cross-pairing, allowing one direction to be resolved along another.

For example, consider
\FieldExampleVectorsDisplay
Both vectors have unit length,
\FieldExampleLengthsDisplay
but their cross-pairing is
\FieldExampleAlignmentDisplay
Since $\FieldExampleVectorA$ is a unit vector, this value is the component of $\FieldExampleVectorB$ along it. Subtracting that component leaves a perpendicular remainder,
\FieldExampleProjectionDisplay
This is the projection operation used below to obtain radial and transverse field directions.

Matrix components admit the same comparison by multiplying corresponding entries and adding. For two real $\FieldExampleMatrixSize$ matrices,
\FieldExampleMatricesDisplay
the pairing is
\FieldExampleMatrixComparisonDisplay
and self-pairing gives
\FieldExampleMatrixLengthDisplay
Transposition and matrix multiplication collect these products on the diagonal of $A^TB$. Their sum, the trace, is therefore the same inner product. Appendix~\ref{app:coupling} shows the entrywise calculation.

To specify the field components, regard a matrix as a linear map on a finite-dimensional real inner-product space $\FieldVectorSpace$. A map is self-adjoint when its action can be transferred between the two entries of every pairing,
\FieldSelfAdjointDisplay
In an orthonormal basis this condition means that its matrix is symmetric. The numerical example and transpose identity in Appendix~\ref{app:coupling} make this equivalence explicit.

Write $\FieldComponentSpace$ for the real vector space of these self-adjoint components. In orthonormal coordinates its inner product is
\FieldMatrixPairingDisplay
The indices range over the basis directions. Positivity follows from the sum of squared entries, while cross-pairing retains the relative alignment needed for the coupling.

\subsection{The complete field and its variations}
Several labelled components require several positions in the field value. For two components, these positions form
\FieldTwoComponentSpaceDisplay
with a typical value
\FieldTwoComponentValueDisplay
Keeping the positions separate allows the same pairing to compare complete field variations by adding the comparisons at matching labels.

A change of both components is recorded as
\FieldTwoComponentChangeDisplay
and its pairing with $\FieldOtherChangePair$ is
\FieldTwoComponentMetricDisplay
This product metric measures the complete change without losing the identity of either component. A worked matrix example is given in Appendix~\ref{app:coupling}.

For a finite label set $\FieldLabelSet$, let $\mathcal V=\FieldComponentSpace^{\FieldLabelSet}$. On a regular configuration domain $\FieldConfigurationDomain$, the field is
\FieldDefinitionDisplay
and its product metric is
\FieldProductMetricDisplay
Coupling and compatibility determine which field values and changes are admissible; the metric compares their magnitudes and directions. Normalizing a component and projecting another perpendicular to it now applies the elementary comparison to the field itself.
